\documentclass[10pt,a4paper]{amsart}
\usepackage[margin=19mm]{geometry}
\usepackage{amsmath,amssymb,mathtools}
\usepackage[scr=rsfs,cal=euler]{mathalfa}
\usepackage{xfrac}
\usepackage[unicode,hidelinks,bookmarks=false]{hyperref}
\newcommand{\ie}{i.e.\ }
\newcommand{\cf}{cf.\ }
\newcommand{\resp}{respectively\ }
\newcommand{\Subsection}{\subsection}
\providecommand{\nobreakdash}{\nobreak\hbox{-}}
\setlength{\emergencystretch}{2em}
\multlinegap0pt
\usepackage{bm}
\usepackage{bbm}
\usepackage{dsfont}
\usepackage{xspace}
\usepackage{cancel}
\usepackage{mathtools}
\usepackage{amsthm}
\makeatletter
\@ifundefined{theorem}{\newtheorem{theorem}{Theorem}}{}
\@ifundefined{assumption}{\newtheorem{assumption}[theorem]{Assumption}}{}
\@ifundefined{lemma}{\newtheorem{lemma}[theorem]{Lemma}}{}
\makeatother
\newcommand{\F}{F}
\newcommand{\R}{{\mathbb R}}
\newcommand{\calX}{{\mathcal X}}
\newcommand{\del}{\delta}
\newcommand{\ds}{\displaystyle}
\newcommand{\iy}{\infty}
\newcommand{\la}{\lambda}
\newcommand{\pl}{\partial}
\title[A Durrett-Remenik particle system in $\mathbb{R}^d$]{A Durrett-Remenik particle system in $\mathbb{R}^d$}
\author{Rami Atar}
\date{Source: arXiv:2507.13990v2 (CC BY 4.0)}
\begin{document}
\maketitle

\noindent\emph{Source and licence.} A Durrett-Remenik particle system in $\mathbb{R}^d$. Version arXiv:2507.13990v2 (CC BY 4.0), distributed under the Creative Commons Attribution 4.0 International licence, as recorded in the arXiv licence metadata for this exact version and shown on the abstract page https://arxiv.org/abs/2507.13990v2. Full rights notice: licenses/arxiv-2507.13990-CC-BY-4.0.txt.\par\smallskip

\noindent\textit{Editorial scope.} Counted unit: the Lemma whose source label is \texttt{lem51} in the article (source0.tex lines 486--489, source label \texttt{lem51}), delivered together with the article's own complete original proof of it (source0.tex lines 491--569, one \texttt{proof} environment of 79 lines). No mathematics is written, repaired or re-derived: statement and proof are the article's own text line for line. Required context retained uncounted: assn10 (lines 358--376); a1 (lines 389--401). Every definition, display and label of those lines is retained unchanged. Omitted from this package and not redistributed: 1--357, 377--388, 402--485, 490--490, 570--1276. Nothing inside a retained excerpt is omitted or altered: every retained body is delivered exactly as the article's lines are written, so no omission receipt of any kind accompanies this package. Citations: the retained excerpts carry no citation command at all, so no bibliography entry of the article is transcribed and no reference is redistributed. Reference adaptation: none was needed and none was made; every reference inside the delivered unit resolves inside it, so no unresolved reference or citation is produced. Printed slips kept literal: no printed word, symbol, hypothesis or number of the counted statement or of its proof was altered. Layout and class: the article's own class is \texttt{amsart} with packages including amsmath, amsfonts, amsthm, amssymb, graphicx, mathtools, amscd, hyperref, eucal, graphics, color, epsfig, cite; the wrapper is the lane's pinned \texttt{amsart} editorial wrapper at 10pt on A4, which restates the theorem-like environments the retained text needs and adds the article's own macro definitions that the retained text needs (\texttt{\textbackslash F}, \texttt{\textbackslash R}, \texttt{\textbackslash calX}, \texttt{\textbackslash del}, \texttt{\textbackslash ds}, \texttt{\textbackslash iy}, \texttt{\textbackslash la}, \texttt{\textbackslash pl}), with extra wrapper packages (bm, bbm, dsfont, xspace, cancel, mathtools). The delivered numbering is the unit's own. Assets: no retained span contains a figure environment, a graphics-inclusion command, a PDF-page inclusion, a table or a TikZ picture; no image file is copied into this package, the package declares no asset dependency and no image file is redistributed with it. Licence: version arXiv:2507.13990v2 of this article is distributed under the Creative Commons Attribution 4.0 International licence, as recorded in the arXiv licence metadata for this exact version and shown on the abstract page https://arxiv.org/abs/2507.13990v2. A full rights notice is delivered at licenses/arxiv-2507.13990-CC-BY-4.0.txt. Characters: every retained excerpt is pure ASCII, so the delivered engine, XeLaTeX with its default Latin Modern text font, reports no missing character.
\begin{assumption}\label{assn10} (i)
$\F\in C(\R^d,\R)$, $\inf_x\F(x)=-\iy$, $\sup_x\F(x)=\infty$, and
for every $a\in\R$, $\F^{-1}\{a\}$ has Lebesgue measure zero.

(ii)
There exists a probability density $\tilde\rho$ and a constant $\tilde c$
such that $\rho(x,y)\le \tilde c\tilde\rho(y-x)$.
Moreover, $\rho(x,y)$ is continuous in $y$ uniformly in $(x,y)$, and
for $a\in\R$ and $x\in \F^{-1}(a,\iy)$, one has
$\int_{\F^{-1}(a,\iy)}\rho(x,y)dy>0$.

(iii)
As $N\to\iy$, $(\bar\xi^N_0,\ell^N_0)\to(\xi_0,\la_0)$ in probability, 
where the latter tuple is deterministic,
$\xi_0(dx)=u_0(x)dx$ and $\la_0\in\R$.
Moreover, $u_0$ is bounded and continuous on $\F^{-1}(\la_0,\iy)$
(and necessarily vanishes on $F^{-1}(-\iy,\la_0)$), and
for every $\del>0$ and $\la\ge \la_0$, $\int_{\F^{-1}(\la,\la+\del)} u_0(x)dx>0$.
\end{assumption}

\begin{equation}\label{a1}
\begin{cases}\ds
{\it (i)}\hspace{3.5em}
\pl_tu(x,t)=\int_{\R^d}u(y,t)\rho(y,x)dy & \quad (x,t):\F(x)>\ell_t,\,t>0,\\
{\it (ii)}\hspace{3.1em}
u(x,t)=0 & \quad (x,t):\F(x)\le\ell_t,\,t>0, \\
{\it (iii)}\hspace{2.5em}
\ds
\int_{\R^d}u(x,t)dx=1, & \quad t>0, \\
{\it (iv)}\hspace{2.9em}
u(\cdot,0)=u_0,\,\ell_0=\la_0.
\end{cases}
\end{equation}

\begin{lemma}\label{lem51}
Let $(u,\ell)\in\calX$ be a solution of \eqref{a1}.
Then $\ell$ is nondecreasing.
\end{lemma}

\begin{proof}
Arguing by contradiction, assume
there exists $t>0$ such that $\ell_t<L_t:=\sup_{s\in[0,t]}\ell_s$.
There are two possibilities.

1. There is $s<t$ such that $\ell_s=L_t$.
In this case, $\ell_\theta\le\ell_s$ for all $\theta\in[s,t]$.
If for some $x$ $\F(x)>\ell_s$ then $\F(x)>\ell_\theta$ for all
$\theta\in[s,t]$, and we can integrate \eqref{a1}(i). Thus
\begin{equation}\label{63}
u(x,t)=u(x,s)+\int_s^t\int_{\R^d}u(y,\theta)\rho(y,x)dyd\theta,
\qquad x\in \F^{-1}(\ell_s,\iy).
\end{equation}
We obtain
\begin{align*}
1=\int_{\R^d} u(x,t)dx&\ge\int_{\F^{-1}(\ell_s,\iy)} \Big[u(x,s)
+\int_s^t\int_{\R^d} u(y,\theta)\rho(y,x)dy d\theta \Big]dx\\
&\ge
1 +\int_s^t\int_{\F^{-1}(\ell_s,\iy)}\int_{\F^{-1}(\ell_s,\iy)} u(y,\theta)\rho(y,x)dy dx d\theta.
\end{align*}
For every $y\in \F^{-1}(\ell_s,\iy)$, one has
$\psi(y):=\int_{\F^{-1}(\ell_s,\iy)}\rho(y,x)dx>0$, by Assumption \ref{assn10}(ii).
Since $\ell_s\ge\ell_\theta$ for all $\theta\in[0,t]$,
we have, similarly to \eqref{63}, that
\[
u(y,\theta)=u_0(y)+\int_0^\theta\int_{\R^d}u(y',\theta')\rho(y',y)dy'd\theta',
\qquad y\in \F^{-1}(\ell_s,\iy).
\]
Hence for such $y$ and all $\theta\in[s,t]$ one has $u(y,\theta)\ge u_0(y)$.
Hence the triple integral above is bounded below by
\[
(t-s)\int_{\F^{-1}(\ell_s,\iy)} u_0(y)\psi(y)dy.
\]
But $\int_{\F^{-1}(\ell_s,\iy)}u_0(y)dy>0$ by assumption, hence
the above integral is positive, a contradiction.

2. There is $s\le t$ such that $\ell_{s-}=L_t$, and $\ell_s<\ell_{s-}$.
In this case there exists $t'>s$ such that $\ell_\theta\le\ell_{s-}$ for all $\theta\in[s,t']$.
We first show
\begin{equation}\label{62}
\int_{\F^{-1}(\ell_{s-},\iy)} u(x,s)dx=1.
\end{equation}
Let $s_n\uparrow s$.
Then $\ell_{s_n}\to\ell_{s-}$ and by assumption,
$\ell_{s_n}\le\ell_{s-}$. Now,
\begin{align}\label{64}
\int_{\F^{-1}(\ell_{s-},\iy)}u(x,s)dx
&=\int_{\F^{-1}(\ell_{s_n},\iy)}u(x,s_n)dx
+\int_{\F^{-1}(\ell_{s_n},\iy)}(u(x,s)-u(x,s_n))dx
\\
&\quad-\int_{\F^{-1}(\ell_{s_n},\ell_{s-})} u(x,s)dx.
\notag
\end{align}
The first term on the right is $1$, and the last term converges to zero
as $n\to\iy$. As for the second term,
since for all $\theta\in[s_n,s]$ one has
$\F^{-1}(\ell_{s-},\iy)\subset \F^{-1}(\ell_{\theta},\iy)$,
one can integrate in \eqref{a1}(i) and get
\[
\int_{\F^{-1}(\ell_{s-},\iy)}|u(x,s)-u(x,s_n)|dx
\le\int_{s_n}^s\int_{\R^d}u(y,\theta)dyd\theta=s-s_n.
\]
The above expression and the second term in \eqref{64}
have the same limit by the assumed boundedness of $u$
on $\R^d\times[0,s]$. This shows \eqref{62}.

Using \eqref{62} and $\F^{-1}(\ell_{s-},\iy)\subset \F^{-1}(\ell_{t'},\iy)$,
we have $\int_{\F^{-1}(\ell_{t'},\iy)} u(x,s)dx=1$.
We can thus repeat the argument above in 1, with
$(s,\ell_{s-},t',\ell_{t'})$ in place of $(s,\ell_s,t,\ell_t)$.
Namely,
\begin{align*}
1=\int_{\F^{-1}(\ell_{t'},\iy)} u(x,t')dx&=\int_{\F^{-1}(\ell_{t'},\iy)} \Big[u(x,s)
+\int_s^{t'}\int_{\R^d} u(y,\theta)\rho(y,x)dy d\theta \Big]dx\\
&\ge
1 +\int_s^{t'}\int_{\F^{-1}(\ell_{s-},\iy)}^\iy\int_{\F^{-1}(\ell_{s-},\iy)} u(y,\theta)\rho(y,x)dy dx d\theta.
\end{align*}
The argument now completes exactly as in case 1.
\end{proof}

\end{document}
